\documentclass[10pt,a4paper]{amsart}
\usepackage[margin=19mm]{geometry}
\usepackage{amsmath,amssymb,mathtools}
\usepackage[scr=rsfs,cal=euler]{mathalfa}
\usepackage{xfrac}
\usepackage[unicode,hidelinks,bookmarks=false]{hyperref}
\newcommand{\ie}{i.e.\ }
\newcommand{\cf}{cf.\ }
\newcommand{\resp}{respectively\ }
\newcommand{\Subsection}{\subsection}
\providecommand{\nobreakdash}{\nobreak\hbox{-}}
\setlength{\emergencystretch}{2em}
\multlinegap0pt
\usepackage{bm}
\usepackage{bbm}
\usepackage{dsfont}
\usepackage{xspace}
\usepackage{cancel}
\usepackage{mathtools}
\usepackage[shortlabels]{enumitem}
\usepackage{amsthm}
\makeatletter
\@ifundefined{theorem}{\newtheorem{theorem}{Theorem}}{}
\@ifundefined{proposition}{\newtheorem{proposition}[theorem]{Proposition}}{}
\makeatother
\newcommand{\R}{\mathbb R}
\newcommand{\X}{\mathcal X}
\newcommand{\longthmtitle}[1]{\mbox{}{\textit{(#1):}}}
\newcommand{\mc}{\mathcal}
\title[Timescale Limits of Linear-Threshold Networks]{Timescale Limits of Linear-Threshold Networks}
\author{William Retnaraj, Simone Betteti, Alexander Davydov, Francesco Bullo, Jorge Cortes}
\date{Source: arXiv:2604.16710v1 (CC BY 4.0)}
\begin{document}
\maketitle

\noindent\emph{Source and licence.} Timescale Limits of Linear-Threshold Networks. Version arXiv:2604.16710v1 (CC BY 4.0), distributed under the Creative Commons Attribution 4.0 International licence, as recorded in the arXiv licence metadata for this exact version and shown on the abstract page https://arxiv.org/abs/2604.16710v1. Full rights notice: licenses/arxiv-2604.16710-CC-BY-4.0.txt.\par\smallskip

\noindent\textit{Editorial scope.} Counted unit: the Proposition whose source label is \texttt{prop:t-LTN\_to\_slow\_limit} in the article (source0.tex lines 517--535, source label \texttt{prop:t-LTN\_to\_slow\_limit}), delivered together with the article's own complete original proof of it (source0.tex lines 537--564, one \texttt{proof} environment of 28 lines). No mathematics is written, repaired or re-derived: statement and proof are the article's own text line for line. Required context retained uncounted: eq:filippov\_reg (lines 235--237); eq:HSS (lines 509--511). Every definition, display and label of those lines is retained unchanged. Omitted from this package and not redistributed: 1--234, 238--508, 512--516, 536--536, 565--832. Nothing inside a retained excerpt is omitted or altered: every retained body is delivered exactly as the article's lines are written, so no omission receipt of any kind accompanies this package. Citations: the retained excerpts carry no citation command at all, so no bibliography entry of the article is transcribed and no reference is redistributed. Reference adaptation: none was needed and none was made; every reference inside the delivered unit resolves inside it, so no unresolved reference or citation is produced. Printed slips kept literal: no printed word, symbol, hypothesis or number of the counted statement or of its proof was altered. Layout and class: the article's own class is \texttt{ieeeconf} with packages including caption, subcaption, xcolor, amsmath, amssymb, upgreek, mathtools, bm, yhmath, mathrsfs, cuted, apptools, algorithmic, enumerate; the wrapper is the lane's pinned \texttt{amsart} editorial wrapper at 10pt on A4, which restates the theorem-like environments the retained text needs and adds the article's own macro definitions that the retained text needs (\texttt{\textbackslash R}, \texttt{\textbackslash X}, \texttt{\textbackslash longthmtitle}, \texttt{\textbackslash mc}), with extra wrapper packages (bm, bbm, dsfont, xspace, cancel, mathtools, enumitem). The delivered numbering is the unit's own. Assets: no retained span contains a figure environment, a graphics-inclusion command, a PDF-page inclusion, a table or a TikZ picture; no image file is copied into this package, the package declares no asset dependency and no image file is redistributed with it. Licence: version arXiv:2604.16710v1 of this article is distributed under the Creative Commons Attribution 4.0 International licence, as recorded in the arXiv licence metadata for this exact version and shown on the abstract page https://arxiv.org/abs/2604.16710v1. A full rights notice is delivered at licenses/arxiv-2604.16710-CC-BY-4.0.txt. Characters: every retained excerpt is pure ASCII, so the delivered engine, XeLaTeX with its default Latin Modern text font, reports no missing character.
\begin{equation}\label{eq:filippov_reg}
    \operatorname{F}[X](x) \coloneqq \bigcap_{\delta>0}\bigcap_{\mu(S)=0} \overline{\operatorname{co}}\{X(B(x, \delta)\setminus S)\}
\end{equation}

\begin{equation}\label{eq:HSS}\tag{HSS}
 x' \in -Dx+\mc H(Ax + u) =: \mc F(x),
\end{equation}

\begin{proposition}\longthmtitle{HSS is the Filippov regularization of the rescaled slow limit}
\label{prop:t-LTN_to_slow_limit}
Let \(F_\infty: \X \rightarrow \R^n\) be defined componentwise as
\begin{equation}\label{eq:slow_pointwise}
(F_\infty(x))_i = \begin{cases}
    - d_ix_i, & A_i^\top x + u_i < 0\\
    0, & A_i^\top x + u_i = 0\\
    1 - d_ix_i, & A_i^\top x + u_i > 0
\end{cases}
\end{equation}
for each \(i \in [n]\).
The following hold:
    \begin{enumerate}[(i)]
        \item For each \(x \in \X\),
        \(\lim_{\tau\to+\infty} F_\tau(x) = F_\infty(x)\);
        \item Assume \(\det A\neq 0\).
        Then, \(\operatorname{F}[F_\infty]\), the Filippov regularization~\eqref{eq:filippov_reg} of \(F_\infty\) is \(\mc F\) as defined in~\eqref{eq:HSS}.
    \end{enumerate}
\end{proposition}

\begin{proof}
    Let \(g(x) = Ax + u\) and let \(g_i(x) = A_i^\top x + u_i\) be its \(i\)th component. 
    For (i), notice that \((F_\tau(x))_i = -d_ix_i + [d_ix_i + \tau g_i(x)]_0^1\).
    If \(g_i(x)>0\), \([d_ix_i + \tau g_i(x)]_0^1 = 1\) for sufficiently large \(\tau\).
    Similarly, if \(g_i(x)<0\), \([d_ix_i + \tau g_i(x)]_0^1 = 0\) for sufficiently large \(\tau\).
    If \(g_i(x) = 0\), and since \(d_ix_i \in [0,1]\), \([d_ix_i + \tau g_i(x)]_0^1 = d_ix_i\).
    Thus, in every case, \(\lim_{\tau\to+\infty} F_\tau(x) = F_\infty(x)\).

    For (ii), we first show that, for each \(x\in\X\), \(\operatorname{F}[F_\infty](x) \subseteq -Dx + \mc H(g(x))\).
    Suppose \(g_i(x) > 0\). 
    By the continuity of \(g_i\) in \(\R^n\), for all \(y\) sufficiently close to \(x\), \((F_\infty(y))_i = 1 - d_iy_i\) and thus \((\operatorname{F}[F_\infty](x))_i = 1 - d_ix_i\).
    Similarly, for \(g_i(x)< 0\), \((\operatorname{F}[F_\infty](x))_i = - d_ix_i\).
    When \(g_i(x) = 0\), we may have points \(y\) such that \(g_i(y)>0\) or \(g_i(y)<0\), so that \((F_\infty(y))_i = 1 - d_iy_i\), \(- d_iy_i\), or \(0\). 
    These are contained in interval \(-d_ix_i + [0,\,1]\), which is their convex closure. 
    Thus, \(\operatorname{F}[F_\infty](x) \subseteq -Dx + \mc H(g(x))\).
    
    We now consider the reverse inclusion.
    Let \(I:=\{i\in[n]:g_i(x)=0\}\). 
    Since \(A\) is invertible, the submatrix \(A_I\) has full row rank.
    Hence for any corner of \(-Dx+\mc H(g(x))\), there exists \(h\in\R^n\) such that \(A_i^\top h>0\) whenever \(\eta_i=1-d_i x_i\) and \(A_i^\top h<0\) whenever \(\eta_i=-d_i x_i\), for all \(i\in I\). 
    Then for \(y_t:=x+th\), one has \(g_i(y_t)=tA_i^\top h\) for \(i\in I\), so its sign agrees with the one corresponding to \(\eta_i\); for \(g_i(x)\neq0\), continuity preserves the sign for all sufficiently small \(t>0\). 
    Thus \(F_\infty(y_t)=\eta\) for all sufficiently small \(t>0\), and in fact on a nonempty open set arbitrarily close to \(x\). 
    This set has positive measure since \(h\) can be perturbed to generate an open set, while still generating the same strict sign in \(A_i^\top \tilde h \lessgtr 0\) for  \(\tilde h\) belonging to that open set, since \(y\mapsto Ay\) is continuous.
    Therefore, removing a null set cannot eliminate all such points, and \(\eta\in \operatorname{F}[F_\infty](x)\). 
    Since \(\operatorname{F}[F_\infty](x)\) is closed and convex (being the intersection of convex closures), it contains the convex hull of all corners of \(-Dx+\mc H(g(x))\), namely \(-Dx+\mc H(g(x))\). 
    Thus \(-Dx+\mc H(g(x))\subseteq \operatorname{F}[F_\infty](x)\).
    Combining with the first inclusion gives \(\operatorname{F}[F_\infty](x)= -Dx+\mc H(g(x))\).
\end{proof}

\end{document}
